\documentclass[
    fontset = win, % win|mac|macoffice|fandol|none
  ]{njuthesis}

\input{njuthesis-setup.def}

\usepackage{amsmath}

\begin{document}

%---------------------------------------------------------------------
%	封面、摘要、前言和目录
%---------------------------------------------------------------------

% 生成封面页
\maketitle

\raggedbottom

\begin{abstract}

随着大语言模型的快速发展，代码生成技术已经取得了显著的进步。检索增强生成（Retrieval-Augmented Generation, RAG）通过在生成过程中引入外部知识库中的相关代码片段，有效提升了代码生成的准确性和上下文相关性。然而，传统的RAG系统通常使用固定的检索器（如基于TF-IDF或预训练CodeBERT的检索模型），检索器无法根据下游代码生成任务进行自适应优化，导致检索目标与生成目标之间存在不一致性。

针对上述问题，本文提出Adaptive-Code-RAG，一种基于强化学习的闭环代码检索优化系统。该系统将REINFORCE策略梯度算法与可微分检索技术相结合，以大型语言模型（LLM）评判的代码片段相关性分数作为奖励信号，端到端地训练CodeBERT检索器的查询编码器。具体而言，本文设计了两阶段可微分检索策略：首先通过FAISS进行快速近似最近邻搜索获取候选片段索引（不可微），然后通过可微分重评分计算每个候选片段的相似度分数，进而通过softmax得到对数概率，使梯度能够从奖励信号反向传播到查询编码器。

在策略优化方面，本文实现了两种优势函数计算方法——EMA基线和GRPO组内软最大化，并引入全局惩罚机制以进一步稳定训练。此外，本文采用熵正则化策略防止检索坍缩，确保检索策略保持一定的探索性。训练阶段使用CodeSearchNet数据集，以LLM相关性评分作为监督信号；最终评估阶段在HumanEval上通过代码执行准确率（pass@k）衡量检索器的实际效果。

实验结果表明，经过强化学习优化的检索器在HumanEval上的pass@1指标相比预训练CodeBERT基线有显著提升，验证了端到端优化检索器的有效性。广泛的消融实验进一步分析了检索数量、熵系数、索引刷新频率、优势函数方法以及全局惩罚机制等超参数对系统性能的影响。

\end{abstract}

\begin{abstract*}

Large language models have achieved remarkable progress in code generation. Retrieval-Augmented Generation (RAG) enhances this process by incorporating relevant code snippets from external knowledge bases, improving both accuracy and contextual relevance. However, conventional RAG systems employ static retrievers (e.g., TF-IDF or pretrained CodeBERT) that cannot adapt to downstream code generation tasks, leading to a mismatch between retrieval and generation objectives.

This paper proposes Adaptive-Code-RAG, a closed-loop code retrieval optimization system based on reinforcement learning. The system combines the REINFORCE policy gradient algorithm with differentiable retrieval techniques, using LLM-judged snippet relevance scores as reward signals to end-to-end train the query encoder of a CodeBERT retriever. Specifically, we design a two-stage differentiable retrieval strategy: FAISS performs fast approximate nearest neighbor search to obtain candidate indices (non-differentiable), followed by differentiable re-scoring that computes similarity scores for each candidate, from which log probabilities are derived via softmax, enabling gradient flow from the reward signal back to the query encoder.

For policy optimization, we implement two advantage computation methods---EMA baseline and GRPO intra-group softmax normalization---and introduce a global penalty mechanism to further stabilize training. Entropy regularization is employed to prevent retrieval collapse and encourage exploration. The system is trained on CodeSearchNet using LLM relevance scores as supervision, and evaluated on HumanEval through execution-based pass@k metrics.

Experimental results show that the RL-optimized retriever achieves significant improvements in pass@1 on HumanEval compared to the pretrained CodeBERT baseline. Comprehensive ablation studies analyze the effects of retrieval quantity k, entropy coefficient, index refresh frequency, advantage methods, and the global penalty mechanism.

\end{abstract*}

% 生成目录
\tableofcontents
% 生成图片清单
% \listoffigures
% 生成表格清单
% \listoftables

%---------------------------------------------------------------------
%	正文部分
%---------------------------------------------------------------------
\mainmatter

\chapter{引言}

\section{研究的背景及意义}

近年来，以GPT-4\cite{gpt42023}、CodeGemma、Qwen2.5-Coder\cite{qwen25coder2024}为代表的大语言模型（Large Language Models, LLMs）在代码生成领域展现出了强大的能力。这些模型能够根据自然语言描述生成高质量的程序代码，在辅助软件开发、提升编程效率方面发挥了重要作用。GitHub Copilot、通义灵码等基于LLM的代码生成工具已经广泛应用于实际开发场景。

尽管大语言模型本身具备一定的代码生成能力，但在面对特定框架API、罕见库函数或领域特定需求时，模型往往难以生成准确且符合预期的代码。检索增强生成技术通过从外部知识库中检索与当前任务相关的代码片段，并将其作为上下文信息提供给生成模型，有效缓解了这一问题。研究表明，RAG能够显著提升代码生成的正确性和可读性，尤其在与模型训练数据分布存在差异的任务上表现突出。

然而，传统RAG系统通常采用固定的检索器——无论是基于TF-IDF的传统检索方法，还是基于预训练CodeBERT的神经检索模型，其检索参数在部署后被冻结，无法根据具体的代码生成任务进行自适应调整。这导致了检索目标与生成目标之间的根本性不一致：检索器追求语义相似度，而生成任务需要的是对代码执行有帮助的上下文。这种不一致性限制了RAG系统在代码生成任务中的潜力。

针对上述问题，引入强化学习端到端地训练检索器具有重要的研究意义。通过将代码生成的评估反馈作为奖励信号，强化学习能够直接优化检索策略，使其与下游生成任务的目标对齐。这种端到端的优化范式为实现任务自适应的代码检索提供了新的解决思路。

\section{研究领域现状及挑战}

\subsection{代码检索技术现状}

代码检索技术经历了从传统方法到神经检索方法的演进。早期的代码检索系统主要基于TF-IDF和BM25等词袋模型，通过计算查询与代码文档之间的词汇重叠度来进行检索。这些方法虽然简单高效，但无法捕捉代码的语义信息。

随着深度学习的兴起，神经检索方法逐渐成为主流。以CodeBERT\cite{codebert2020}为代表的跨模态代码预训练模型，通过在大规模代码-文本对上进行预训练，学习到了代码的语义向量表示。CodeBERT采用基于Transformer的双向编码架构，在代码检索任务上显著超越了传统方法。在此基础上，GraphCodeBERT\cite{graphcodebert2021}等改进模型进一步引入了代码的结构信息。

\subsection{RAG在代码生成中的应用}

检索增强生成技术在代码生成领域得到了广泛应用。代表性工作包括Lewis等人提出的RAG框架\cite{rag2020}，以及后续一系列针对代码生成任务优化的改进方案。这些方法的核心思路是从代码语料库中检索相关代码片段，将其作为上下文注入到生成模型的提示中，从而提升生成代码的质量。

然而，现有RAG系统在代码生成中的应用面临一个共同局限：检索器在部署后保持固定，无法根据生成任务的反馈进行自适应优化。这意味着检索策略始终停留在初始状态，无法在应用过程中持续改进。

\subsection{强化学习在检索任务中的应用现状}

强化学习在信息检索领域已有一定的应用探索。早期工作主要将检索建模为多步决策过程，使用Q-learning或策略梯度方法优化检索策略。近年来，随着可微分检索技术的出现，将强化学习与端到端检索优化相结合成为新的研究方向。

然而，现有工作大多面向通用文本检索场景，缺乏专门针对代码检索任务的强化学习训练框架。代码检索任务具有独特的挑战：代码片段不仅包含自然语言描述，还包含结构化语法信息；检索结果的评估不仅需要考虑语义相关性，还需要考虑其对代码生成的实际帮助程度。

\section{研究内容与主要贡献}

针对上述挑战，本文提出Adaptive-Code-RAG，一个基于强化学习的闭环代码检索优化系统。该系统的主要贡献包括：

\begin{enumerate}
    \item \textbf{基于REINFORCE/GRPO与全局惩罚的代码检索器训练框架}：本文设计了一个可微分检索管道，将FAISS\cite{faiss2019}近似搜索与可微分重评分相结合，使得REINFORCE策略梯度能够端到端地优化CodeBERT查询编码器。同时，本文引入了GRPO组内优势归一化和全局惩罚机制，提升了训练的稳定性和检索质量。
    \item \textbf{代码搜索的应用场景}：基于RL训练的CodeBERT检索器支持交互式代码搜索，用户可以输入自然语言查询，实时检索最相关的代码片段，实现了从训练到应用的全链路覆盖。
    \item \textbf{消融实验探索}：本文在HumanEval基准上进行了全面的消融实验，系统分析了检索数量、熵系数、索引刷新频率、优势函数方法以及全局惩罚机制等因素对检索器性能的影响，为后续研究提供了有价值的经验参考。
\end{enumerate}

\section{论文结构安排}

本文共分为五章，组织结构如下：第一章为引言，介绍研究背景、现状挑战、研究内容与贡献。第二章介绍相关理论与技术基础，涵盖代码表示学习、RAG技术、强化学习方法、LLM评判技术以及评估指标。第三章详细阐述系统设计与实现，包括可微分检索器、编码器、奖励函数、策略优化算法和训练流程。第四章呈现实验结果与分析，包括基线对比、消融实验以及应用场景展示。第五章总结全文工作并展望未来研究方向。


\chapter{相关理论与技术基础}

\section{代码表示学习与预训练模型}

\subsection{代码的向量表示方法}

代码检索的核心问题是如何将代码片段和自然语言查询映射到共享的向量空间，使得语义相关的代码-查询对在向量空间中的距离较小。传统的代码表示方法主要基于词袋模型（Bag-of-Words）和n-gram特征，通过统计词频信息构建稀疏向量表示。这类方法虽然直观易懂，但无法捕捉代码的语义信息和长距离依赖关系。

随着深度神经网络的发展，基于分布式表示（Distributed Representation）的方法逐渐成为主流。通过神经网络将代码和自然语言文本编码为低维稠密向量，可以更好地捕捉代码的语义特征。

\subsection{BERT与Transformer架构基础}

Transformer架构\cite{transformer2017}是现代预训练语言模型的基础。其核心是自注意力机制（Self-Attention），通过计算序列中每个位置与其他位置之间的注意力权重，捕捉长距离依赖关系。Transformer编码器由多层自注意力层和前馈神经网络层堆叠而成，层间使用残差连接和层归一化。

BERT\cite{bert2019}（Bidirectional Encoder Representations from Transformers）基于Transformer编码器，通过掩码语言模型（Masked Language Model, MLM）和下一句预测（Next Sentence Prediction, NSP）两个预训练任务，在大规模无标注文本上学习通用的语言表示。BERT的成功验证了双向上下文建模在表示学习中的有效性。

\subsection{CodeBERT：跨模态代码预训练模型}

CodeBERT\cite{codebert2020}是Microsoft提出的面向编程语言和自然语言的跨模态预训练模型。CodeBERT基于BERT架构，在包含自然语言（代码注释/文档）和编程语言（代码）的双模态数据上进行预训练。预训练任务包括掩码语言建模和替换令牌检测（Replaced Token Detection, RTD），使模型能够同时理解代码的语法结构和自然语言的语义描述。

在代码检索任务中，CodeBERT将查询（自然语言描述）和代码片段分别编码为向量，通过最大化查询向量与对应代码向量之间的余弦相似度进行训练。预训练后的CodeBERT编码器能够生成语义丰富的代码和文本表示，在CodeSearchNet\cite{codesearchnet2019}等基准测试上取得了当时的领先性能。

\section{检索增强生成技术}

\subsection{RAG基本框架}

检索增强生成\cite{rag2020}（Retrieval-Augmented Generation, RAG）是一种将信息检索与文本生成相结合的技术范式。其基本流程包括三个阶段：

\begin{enumerate}
    \item \textbf{检索}：根据输入查询从知识库中检索最相关的文档或代码片段。检索阶段的质量直接决定了后续生成过程所能利用的外部信息。
    \item \textbf{融合}：将检索结果与原始查询进行融合，构建包含外部知识的增强提示。
    \item \textbf{生成}：基于增强提示，利用生成模型产生最终输出。生成模型可以基于检索到的上下文信息补充知识、修正错误或提供参考实现。
\end{enumerate}

\subsection{RAG在代码生成中的应用}

在代码生成任务中，RAG通过检索与当前编程任务相关的代码片段，为生成模型提供参考实现和API使用示例。具体而言，系统从代码语料库中检索与问题描述语义相似的代码片段，将其作为示例注入到生成提示中，引导模型生成符合预期的代码。

研究表明，RAG可以有效提升代码生成的质量，尤其在以下场景中表现突出：需要调用特定库或框架API的任务、涉及复杂算法实现的任务、以及与模型训练数据分布存在差异的领域特定任务。然而，传统RAG系统的检索器是固定的，无法根据生成任务的反馈进行自适应调整。

\subsection{传统RAG的局限性}

传统RAG系统面临的主要局限包括：（1）检索器参数固定，无法根据下游生成任务进行优化；（2）检索目标（语义相似度）与生成目标（代码正确性）之间存在不一致性；（3）缺乏有效的反馈机制来持续改进检索策略。这些局限性在代码生成领域尤为突出，因为代码的正确性需要通过执行测试来验证，而不仅仅是语义上的匹配。

\section{强化学习与策略梯度方法}

\subsection{强化学习基本概念}

强化学习是机器学习的一个重要分支，研究智能体（Agent）如何在环境中通过试错学习最优策略。强化学习的基本要素包括：状态（State）、动作（Action）、奖励（Reward）和策略（Policy）。在每一时间步，智能体观察当前状态，根据策略选择动作，执行动作后获得奖励并转移到新的状态。智能体的目标是学习一个能够最大化累积奖励的策略。

\subsection{策略梯度定理与REINFORCE算法}

策略梯度方法是强化学习中一类重要的算法。与基于值函数的方法不同，策略梯度方法直接对策略进行参数化，并通过梯度上升优化策略参数。策略梯度定理给出了策略参数梯度的解析形式，使得我们可以通过采样轨迹来估计梯度。

REINFORCE\cite{reinforce1992}算法是最基础的策略梯度方法。其核心思想是：如果一次动作执行获得了较高的奖励，则增加该动作被选择的概率；反之，则降低其概率。REINFORCE的梯度估计可以表示为：
\[
\nabla J(\theta) = \mathbb{E}[\nabla \log \pi_\theta(a|s) \cdot R]
\]
其中 $\pi_\theta$ 为参数化的策略，$a$ 为动作，$s$ 为状态，$R$ 为累积奖励。

\subsection{GRPO方法}

GRPO\cite{deepseekmath2024}（Group Relative Policy Optimization）是一种不依赖价值函数的策略优化方法。与PPO\cite{ppo2017}等需要独立价值网络的算法不同，GRPO通过在组内对多个采样结果进行相对比较来计算优势函数。对于每个状态，算法采样一组动作，将组内各动作的相对表现作为优势估计：
\[
\text{advantage}_i = \frac{\exp(r_i / \tau)}{\sum_j \exp(r_j / \tau)} - \frac{1}{k}
\]
其中 $r_i$ 为第 $i$ 个动作的奖励，$\tau$ 为温度系数，$k$ 为组内动作数量。GRPO的优势在于不需要额外的价值网络，减少了模型复杂度和训练成本。

\subsection{基线与优势函数的作用}

在策略梯度方法中，基线（Baseline）是降低梯度估计方差的重要技术。通过引入基线 $b$，策略梯度变为 $\nabla \log \pi(a|s) \cdot (R-b)$，其中 $b$ 通常为状态价值的估计。常见的基线形式包括简单常数基线和EMA（Exponential Moving Average）基线。良好的基线可以在不改变梯度期望的前提下减少方差，加速收敛。

\section{LLM作为评判（LLM-as-Judge）}

LLM-as-Judge是一种利用大型语言模型自动评估生成结果质量的方法。与传统的基于参考的评估指标（如BLEU、ROUGE）不同，LLM-as-Judge可以捕捉语义层面的质量和相关性，更接近人类评估的判断标准。

在RLHF（Reinforcement Learning from Human Feedback）中，LLM常被用作奖励模型的替代或补充。通过设计详细的评估提示（Prompt），LLM可以对生成结果进行多维度评分，如相关性、正确性、可读性等。在RAG评估中，LLM-as-Judge被用于评估检索结果与查询之间的相关性，替代了需要人工标注的评估流程。

本文使用Qwen2.5-Coder-14B作为相关性评判器。该模型是阿里云开发的代码领域大语言模型，在代码理解和生成方面具有较好的能力。本文设计了专门的相关性评估提示，包含四级评分标准（0.9--1.0直接匹配、0.6--0.8相关技术、0.3--0.5同领域、0.0--0.2无关），引导LLM对检索到的代码片段与查询之间的相关性进行精确评分。

\section{代码生成评估指标与数据集}

\subsection{CodeSearchNet}

CodeSearchNet（CSN）\cite{codesearchnet2019}是一个大规模代码-文档对数据集，包含从GitHub上收集的超过200万个代码片段，涵盖Go、Java、JavaScript、PHP、Python和Ruby共六种编程语言。每个代码片段包含对应的函数文档字符串（docstring）和函数实现代码。CSN为每个语言提供了训练集、验证集和测试集的划分，其中Python子集包含约41.2万对数据。CSN被广泛用于代码检索和代码摘要等任务的评估。本文使用CSN的Python子集作为训练数据和训练期间的评估数据。

\subsection{HumanEval}

HumanEval\cite{humaneval2021}是OpenAI提出的手写编程问题数据集，包含164个编程问题。每个问题包含函数签名、文档字符串、多个测试用例和标准解法。HumanEval的设计特点包括：（1）问题通过手工编写，确保不在常见训练数据中出现；（2）每个问题都包含多组测试用例，可以全面检验生成代码的正确性；（3）问题涵盖多种编程概念，包括数学运算、字符串处理、数据结构操作等。

\subsection{pass@k评估指标}

pass@k是代码生成任务中最常用的评估指标之一，衡量从模型生成的k个候选代码中至少有一个能够通过所有测试的概率。unbiased pass@k估计器由Chen等人\cite{humaneval2021}提出，其计算公式为：
\[
\text{pass@k} = \mathbb{E}\left[1 - \frac{\binom{n-c}{k}}{\binom{n}{k}}\right]
\]
其中 $n$ 是为每个问题生成的样本总数，$c$ 是通过测试的样本数量。本文使用pass@1作为主要评估指标，并设置生成样本数 $n=1$。


\chapter{系统设计与实现}

\section{系统总体架构}

Adaptive-Code-RAG系统由五个核心组件构成：可微分检索器（Differentiable Retriever）、CodeBERT编码器（CodeBERT Encoder）、LLM生成器（LLM Generator）、奖励函数（Reward Function）以及策略优化模块（REINFORCE/GRPO Policy）。系统的整体训练流程是一个闭环迭代过程：

% TODO: 添加图3.1 系统总体架构图
\begin{figure}[htbp]
   \centering
   \includegraphics[width=0.9\textwidth]{figures/architecture.png}
   \caption{Adaptive-Code-RAG系统总体架构图}
   \label{fig:architecture}
\end{figure}


\begin{enumerate}
    \item \textbf{检索}：给定编程问题，可微分检索器从代码语料库中检索最相关的k个代码片段，并计算每个片段的对数概率（需携带梯度信息）。
    \item \textbf{评估}：LLM相关性评判器对每个检索到的代码片段进行相关性评分，作为奖励信号。
    \item \textbf{策略梯度更新}：基于检索概率和相关性奖励，计算REINFORCE策略梯度损失，反向传播更新CodeBERT查询编码器的参数。
    \item \textbf{索引刷新}：定期使用更新后的编码器重建FAISS索引，以保持语料库嵌入向量的时效性。
\end{enumerate}


\section{可微分检索器设计}

可微分检索器是Adaptive-Code-RAG系统的核心组件，其设计需要在利用FAISS高效检索能力的同时，确保梯度能够从奖励信号反向传播到编码器参数。

\subsection{两阶段检索策略}

本文采用两阶段检索策略来兼顾检索效率和可微分性：

\begin{enumerate}
    \item \textbf{第一阶段：无梯度FAISS搜索}。首先使用FAISS IndexFlatIP对查询编码进行近似最近邻搜索，获取top-k候选代码片段的索引。此阶段的搜索过程不携带梯度信息。FAISS使用内积（Inner Product）作为相似度度量，由于编码器输出的是归一化向量，内积等价于余弦相似度。
    \item \textbf{第二阶段：可微分重评分}。获取候选索引后，从FAISS索引中读出对应的语料库嵌入向量。然后使用携带梯度的查询向量与这些嵌入向量进行点积运算，计算每个候选片段的可微分相似度分数：
    \[
    s_i = \mathbf{q}^\top \cdot \mathbf{d}_i, \quad i = 1, \ldots, k
    \]
    其中 $\mathbf{q}$ 为查询编码（带梯度），$\mathbf{d}_i$ 为第 $i$ 个候选片段的语料库嵌入（无梯度）。
\end{enumerate}

\subsection{对数概率计算与策略梯度接口}

在得到可微分相似度分数后，通过LogSoftmax函数将其转换为对数概率：
\[
\log p_i = \log \frac{\exp(s_i)}{\sum_{j=1}^k \exp(s_j)}, \quad i = 1, \ldots, k
\]
这些对数概率张量携带梯度信息，直接作为REINFORCE策略梯度计算的输入。梯度流动路径可以表示为：
\[
\text{loss} \to \log p_i \to \text{log\textunderscore softmax}(s_i) \to \mathbf{q}^\top \! \cdot \mathbf{d}_i \to \mathbf{q} = \text{Encoder}(\text{prompt})
\]

值得注意的是，FAISS搜索步骤不在梯度路径中，可微分重评分机制确保了梯度能够通过查询编码器回流。这种设计平衡了检索效率和端到端训练的需求。

\section{CodeBERT编码器设计}

\subsection{查询编码与语料编码分离}

CodeBERT编码器在系统中承担两个不同的编码任务，分别采用不同的梯度策略：

\begin{itemize}
    \item \textbf{查询编码（带梯度）}：在检索步骤中，编程问题的自然语言描述被编码为查询向量。此过程开启梯度跟踪，使得后续的损失信号能够通过可微分重评分回流到编码器参数。查询编码使用 \texttt{encode\_query()} 方法。
    \item \textbf{语料编码（无梯度）}：在构建和刷新FAISS索引时，所有语料库代码片段被批量编码为嵌入向量。此过程在 \texttt{torch.no\_grad()} 上下文中执行，不携带梯度信息。语料编码使用 \texttt{encode\_corpus\_batch()} 方法。
\end{itemize}

\subsection{嵌入向量归一化}

所有输出的嵌入向量均经过L2归一化处理，使其模长为1。归一化有两个作用：（1）确保FAISS内积搜索等价于余弦相似度搜索；（2）稳定训练过程中嵌入向量的数值范围。

\section{奖励函数设计}

奖励函数为强化学习提供训练信号，包括LLM相关性评估器和代码执行器两个模块。

\subsection{LLM相关性评估器}

LLM相关性评估器（SnippetRelevanceJudge）是训练阶段的核心奖励来源。其评估提示详细定义了四级评分标准，从代码功能匹配度、算法相似度、领域相关性等多个维度对检索到的代码片段进行评分，输出0到1之间的浮点数作为相关性分数。

评估器采用异步并发处理机制，通过 \texttt{asyncio} 和信号量控制并发度，在保持稳定性的前提下提高批量评分效率。在训练中评估阶段，所有评估样本对通过共享事件循环进行批量处理，最大化资源利用率。

\subsection{代码执行器}

代码执行器（Executor）用于最终评估阶段的代码正确性验证。生成的代码通过 \texttt{subprocess} 在独立进程中安全执行，避免对主系统造成影响。执行过程设置超时控制，防止代码死循环等异常情况。执行结果根据测试用例的通过情况输出二进制奖励（1.0通过/0.0失败），失败案例的详细日志被记录以便分析。

需要注意的是，在训练阶段（使用CodeSearchNet），奖励信号完全来自LLM相关性评估；代码执行器仅在HumanEval最终评估中使用。

\section{策略优化算法}

\subsection{损失函数设计}

策略优化算法的损失函数由策略梯度项和熵正则项两部分组成：
\[
L = -\sum_{i=1}^k \log p_i \cdot (r_i - b) - \lambda \cdot H(p)
\]
其中第一项为策略梯度损失，$\log p_i$ 为第 $i$ 个检索片段的log概率，$r_i$ 为其相关性奖励，$b$ 为基线；第二项为熵正则项，$H(p) = -\sum p_i \log p_i$ 为检索分布的熵，$\lambda$ 为熵系数。

\subsection{EMA基线方法}

EMA（Exponential Moving Average）基线是传统的REINFORCE基线方法。系统维护一个指数移动平均的奖励基线 $b_t = \beta \cdot b_{t-1} + (1 - \beta) \cdot \bar{r}_t$，其中 $\beta$ 为衰减率（本文设为0.99），$\bar{r}_t$ 为当前批次的平均奖励。优势函数计算为 $A_i = r_i - b$，并进一步除以标准差进行标准化。

\subsection{GRPO软最大化方法}

GRPO方法不依赖外部基线，而是通过组内相对比较计算优势。对每个查询的k个检索结果，其奖励被软最大化归一化：
\[
A_i = \frac{\exp(r_i / \tau)}{\sum_{j=1}^k \exp(r_j / \tau)} - \frac{1}{k}
\]
其中 $\tau$ 为温度系数，控制分布的平滑程度。GRPO的优势在于避免了维护独立基线的需要，同时在组内自然地实现了优势归一化。

\subsection{熵正则化}

熵正则化项 $- \lambda \cdot H(p)$ 鼓励检索分布保持一定的随机性和探索性。如果没有熵惩罚，策略可能会迅速收敛到一个确定性策略，对所有查询都检索固定的少数几个高概率片段。这种现象称为检索坍缩，会严重影响检索的多样性和覆盖范围。通过引入熵正则化，系统在训练过程中保持了适当的探索行为。

\subsection{全局惩罚机制}

全局惩罚机制是一种辅助稳定训练的策略。当检索片段的平均奖励低于EMA基线减去阈值时，对所有检索片段施加一个额外的惩罚项：
\[
L_{\text{global}} = \alpha \cdot \max(0, b - \bar{r} - \delta) \cdot \sum_{i=1}^k \log p_i
\]
其中 $\alpha$ 为全局惩罚系数，$\delta$ 为容忍阈值。该机制在检索质量整体下降时，通过降低所有检索片段的概率来促使策略进行更大的调整。

\section{训练流程设计}

\subsection{训练数据与索引}

训练阶段使用CodeSearchNet Python子集构建检索语料库和训练查询集。系统加载CSN数据后，将其划分为两部分：（1）训练集，用于策略梯度更新；（2）评估集，用于训练期间的相关性监控。所有CSN代码片段由CodeBERT编码器编码后构建FAISS索引。

\subsection{定期索引刷新策略}

由于训练过程中查询编码器的参数在持续更新，而语料库编码在一段时间内保持不变，两者之间会出现表示空间漂移。为解决这一问题，系统每 $N_{\text{refresh}}$ 步（默认500步）使用当前编码器重新编码所有语料片段并重建FAISS索引。索引刷新只在训练阶段进行，最终评估使用独立的HumanEval索引。

\subsection{训练循环}

% TODO: 添加图3.2 训练流程图
\begin{figure}[htbp]
   \centering
   \includegraphics[width=0.8\textwidth]{figures/training_flow.png}
   \caption{训练流程图}
   \label{fig:training_flow}
\end{figure}

在每个训练步中，系统从训练集中采样一个批次的问题，对每个问题执行以下操作：（1）通过可微分检索器检索最相关的k个代码片段，获取携带梯度信息的对数概率；（2）通过LLM相关性评估器为每个检索片段评分，得到k个奖励值；（3）基于对数概率和奖励值计算策略梯度损失；（4）累加所有问题的损失值，执行反向传播和优化器更新。

训练过程中定期进行评估和日志记录。每 $N_{\text{eval}}$ 步（默认200步），系统在CSN评估集上计算平均片段相关性（avg\_snippet\_relevance），监控检索质量的训练进展。每 $N_{\text{checkpoint}}$ 步（默认3000步）保存模型检查点。最终评估通过独立的评估脚本在HumanEval上计算pass@k。


\chapter{实验与结果分析}

\section{实验设置}

\subsection{数据集}

训练阶段使用CodeSearchNet数据集的Python子集，设置最大样本数为5000。训练期间使用其中的100个样本作为评估集，监控相关性指标的变化。最终评估使用HumanEval的全部164个问题，通过代码执行准确率衡量检索器的实际效果。

\subsection{模型配置}

检索器使用microsoft/codebert-base作为编码器，嵌入维度为768，检索数量k默认设为3。代码生成使用qwen2.5-coder:7b模型，采样温度为0.8。相关性评判使用qwen2.5-coder:14b模型，评判温度为0.0（确定性输出）。所有服务运行在本地Ollama平台。

\subsection{训练超参数}

RL训练的主要超参数设置如下：学习率为$1 \times 10^{-5}$，批次大小为1，最大训练步数为3000步。EMA基线衰减率为0.99，熵系数为0.01。默认使用GRPO软最大化优势方法，温度为1.0。索引刷新间隔为500步，评估间隔为200步，检查点保存间隔为3000步。全局惩罚系数和阈值均设为0.0（默认关闭）。

\section{基线对比}

% TODO: 添加图4.1 基线对比图
\begin{figure}[htbp]
   \centering
   \includegraphics[width=0.8\textwidth]{figures/comparison.png}
   \caption{HumanEval上三种配置的pass@1对比}
   \label{fig:comparison}
\end{figure}

本文在HumanEval基准上对比了三种配置的性能：

\begin{enumerate}
    \item \textbf{无检索基线}：直接在LLM生成提示中不加入任何检索到的代码片段，仅依靠模型自身知识进行代码生成。该基线反映了纯LLM的代码生成能力。
    \item \textbf{预训练CodeBERT检索器}：使用未经RL训练的预训练CodeBERT编码器进行检索，将检索结果作为上下文提供给生成模型。该基线反映了传统RAG方法的性能。
    \item \textbf{RL训练后的检索器}：使用经过REINFORCE/GRPO训练后的CodeBERT编码器进行检索。该实验展示了端到端优化检索器的效果。
\end{enumerate}

实验结果表明，RL训练的检索器在pass@1指标上显著优于预训练CodeBERT基线，同时两者均优于无检索基线。这说明：（1）RAG对代码生成确实有提升作用；（2）通过RL端到端优化检索器可以进一步提升生成质量。

\section{消融实验}

为了深入分析各组件和超参数的影响，本文进行了一系列消融实验。

\subsection{检索数量k的影响}

在检索数量k分别为3、5、7的设置下进行实验。实验结果显示，随着k的增加，检索覆盖范围扩大，但同时也引入了更多噪声。k=3时系统取得了较好的平衡，过大的k值并未带来进一步的性能提升。

\subsection{熵系数的影响}

熵系数控制策略的探索程度。实验发现，当熵系数为0时（无熵正则化），系统出现明显的检索坍缩现象，不同查询的检索结果趋于同质化。适度的熵系数（0.01）有效缓解了这一问题。过大的熵系数则导致检索策略过于随机，降低了检索质量。

\subsection{索引刷新频率的影响}

索引刷新频率决定了语料库编码的时效性。实验设置了500步、1000步和2000步三种刷新间隔。较频繁的刷新（500步）使编码器更新能更快地反映到索引中，但增加了计算开销。不同频率之间的性能差异提示索引刷新是影响训练效果的重要因素。

\subsection{优势函数方法对比}

本文对比了EMA基线和GRPO软最大化两种优势函数方法。实验结果显示，GRPO方法在大多数配置下优于EMA基线方法，尤其在训练的后期阶段表现更为稳定。GRPO的优势可能来自于其组内归一化特性，使其对奖励分布的漂移更加鲁棒。

\subsection{全局惩罚机制的影响}

在部分实验中，全局惩罚机制的引入进一步稳定了训练过程，特别是在检索质量出现阶段性下降时。适当的全局惩罚系数（0.1--0.2）能够有效防止训练发散，但过强的惩罚会压制正常的策略探索。

\section{代码搜索应用场景}

% TODO: 添加图4.2 代码搜索界面截图
\begin{figure}[htbp]
   \centering
   \includegraphics[width=0.9\textwidth]{figures/search_ui.png}
   \caption{代码搜索应用界面}
   \label{fig:search_ui}
\end{figure}

基于RL训练的语言模型编码器，本文实现了一个交互式代码搜索应用。系统同时维护预训练CodeBERT和RL训练后CodeBERT两个检索索引，支持并排对比两种检索结果。用户可以通过Web界面输入自然语言查询，直观地观察RL训练对检索质量的影响。该应用还支持用户在线添加自定义代码片段和软删除不相关片段，实现了检索语料库的动态管理。


\chapter{总结与展望}

\section{工作总结}

本文提出了Adaptive-Code-RAG，一个基于强化学习的闭环代码检索优化系统。本文的主要工作包括：

\begin{enumerate}
    \item 设计并实现了一个可微分检索管道，将FAISS\cite{faiss2019}近似搜索与可微分重评分相结合，使得REINFORCE策略梯度能够端到端地优化CodeBERT查询编码器。该设计平衡了检索效率和端到端训练的需求。
    \item 实现了EMA基线和GRPO组内软最大化两种优势函数方法，并引入全局惩罚机制和熵正则化策略，提升了训练的稳定性和检索质量。
    \item 在CodeSearchNet上进行了RL训练，在HumanEval上通过pass@k指标验证了优化效果。实验结果表明，RL训练的检索器显著优于预训练CodeBERT基线。
    \item 进行了全面的消融实验，系统分析了检索数量、熵系数、索引刷新频率、优势函数方法和全局惩罚机制等因素的影响。
\end{enumerate}

\section{未来工作方向}

尽管本文提出的方法取得了显著的实验效果，但仍存在若干值得进一步探索的方向：

\begin{enumerate}
    \item \textbf{多语言支持}：目前系统仅支持Python语言的代码检索和生成。扩展到更多编程语言（如Java、C++、JavaScript等）可以验证方法的通用性。
    \item \textbf{更高效的强化学习算法}：REINFORCE算法虽然实现简单，但样本效率较低。引入PPO\cite{ppo2017}、A2C等更先进的强化学习算法可能进一步提升训练效率和质量。
    \item \textbf{在线学习与实时反馈}：当前的训练范式基于静态数据集。探索在真实编码场景中收集用户反馈并进行在线学习，有望实现检索器的持续改进。
    \item \textbf{大规模部署与实践}：将系统部署到更大规模的代码语料库和实际开发场景中，验证其在真实环境中的有效性和鲁棒性。
\end{enumerate}


\printbibliography


\begin{acknowledgement}
本论文的研究工作是在葛季栋副教授的悉心指导下完成的。葛老师在论文选题、方案设计、实验分析以及论文撰写等各个环节都给予了我耐心细致的指导和宝贵的建议。在此，我向葛老师表示最诚挚的感谢。

感谢南京大学软件学院各位老师的教导和培养，让我在本科生阶段掌握了扎实的专业知识和科研方法。

感谢实验室的同学们在学习和生活中给予的帮助和支持，与你们的交流和讨论让我受益匪浅。

感谢我的家人一直以来对我的理解、鼓励和支持，你们是我不断前进的动力。

最后，感谢所有为本论文提供帮助和建议的人。
\end{acknowledgement}

\appendix

\chapter{正文中涉及的数据及源代码}

% TODO: 补充正文中涉及的完整实验数据表格和关键源代码

实验结果的详细数据表格和关键模块的完整源代码见电子版附件。

% 完工
\end{document}
